While this paper represents an initial step in the development of observability metrics and the exploration of optimization strategies, it is important to acknowledge certain limitations. 
First, the approach relies on simulation and isolated experiments, which may not fully capture the complexities and faults faced in real-world applications. For this purpose, we deliberately designed \name{} with extensibility in mind, allowing developers to integrate their own load curves and custom faults. In the future, we aim to further validate the approach under real loads and more extensive fault scenarios.  

Second, the individual metrics proposed for fault observability were deliberately kept clear and flexible, to enable a technology-independent assessment. A generalization to more complex fault detection systems might not be possible without high coupling to the reporting mechanism.  Still, we consider our approach validated as we were able to show the impact of different configuration options on the generated observability data, evident both in the graphical plots and in the classifier score. As part of our ongoing effort to expand observability assessments, we aim to enhance our tool by incorporating support for alerting systems and Mean Time To Detect (MTTD). Besides this, we also aim to improve trade-off analysis by incorporating other cost metrics beyond CPU utilization.

Third, the current implementation of \name{} has some limitations, notably the absence of certain features like support for logs\footnote{At the time of implementation, \textsc{OpenTelemetry}  had not finalized logs.}, \textsc{Kubernetes} integration, among others. We aim to address these as we continue to utilize our tool in further observability assessments.
In future work, we plan to explore how observability decisions can be optimized and automated. Taking our metrics as a foundation, our next objective is to delve into intelligent and learning approaches to optimize and tune observability configuration parameters. 







